\exnumber{05}
\extitle{UX V - Ex 6 - Protocole de test}
\exfiles{ex05/PDF, nomme Name\_FirstName\_DeliverableName\_Date\_VersionNumber}
\exauthorize{Ceux nommés dans l'exercice}
\makeheaderfiles

\textbf{Instructions :}

Préparez un protocole pour des tests d'utilisabilité.

\textbf{Livrable :}

\begin{itemize}
\item Un document PDF avec un protocole de test
\item Déposez le PDF sur Github, nommé ainsi : « Name\_FirstName\_DeliverableName\_Date\_VersionNumber ». Exemple : « Doe\_John\_MidFiPrototype\_05242024\_V1 ».
\item Ajoutez un titre dans le document.
\item Soignez l'aspect global du document. Il doit être clair, simple et lisible. Vous pouvez regarder canva.com, sans surcharger le dessin de détails.
\end{itemize}

\textbf{Pistes :} Étapes recommandées

\begin{itemize}
\item Définir le périmètre du test
\item Recruter 5 utilisateurs
\item Identifier les objectifs : que prouvez-vous ?
\item Établir des métriques : comment prouvez-vous que votre design est efficace ?
\item Décrire les étapes et les consignes
\item Rédiger une introduction pour accueillir l'utilisateur
\item Rédiger une conclusion pour remercier l'utilisateur
\end{itemize}

\textbf{Ressources}

\begin{itemize}
\item Comment mener un entretien utilisateur (tutoriel de Google Ventures) \url{https://www.youtube.com/watch?v=Qq3OiHQ-HCU}
\item Mener un test d'utilisabilité \url{https://www.usability.gov/how-to-and-tools/methods/running-usability-tests.html}
\item Tests d'utilisabilité : qu'est-ce que c'est et comment les faire ? \url{https://uxdesign.cc/usability-testing-what-is-it-how-to-do-it-51356e5de5d}
\end{itemize}

\textbf{Citation :}

Vous devez détecter : Combien d'erreurs les utilisateurs font-ils ? Quelle est la gravité de ces erreurs ? À quel point peuvent-ils facilement se rétablir des erreurs ?

